\section{A smaller host with one literal cyclic character}
\label{sec:cyclic-cm-host}

The noncyclicity in AC-G30 belongs to the chosen signed
Jacobian host.  It is not an obstruction to every bare Host
containing the cubic coefficient field.  The following smaller
construction removes that particular field-definition gate.
It retains an explicit integral CM-type gate.

\begin{proposition}[AC-G32: a cubic-compatible cyclic Host]
\label{AC-G32}
Let $h\geq64$ and put
\[
 R_h=\mathbb Q(\zeta_{2^{h+2}})^+,
 \qquad K=\mathbb Q(\omega),\qquad C_h=R_hK,
 \qquad n=2^{h+1},\qquad Q_S=3\cdot2^{h+2}.
                                                               \tag{G32.1}
\]
Use the frame
$A=\operatorname{Gal}(C_h/\mathbb Q)\simeq C_{2^h}\times C_2$,
projecting to its first factor and with distinguished generator
induced by $5$.  Then $n_0=2^h$ and $o=2$.
Take any source \texttt{Ranks} for this frame.  For example
$k=2n$, $m=2$, $t=k+7$ are permissible.
Let $B\geq269$, and choose distinct primes
$q,d>\max(B,4t(k+m-1)+1,11)$ satisfying
\[
 \begin{gathered}
 q\equiv3\pmod8,\quad d\equiv7\pmod {16},\quad
 d\equiv-1\pmod q,\\
 d\equiv-1\pmod\ell
       \quad\text{for every odd prime }\ell\leq B.
 \end{gathered}
                                                               \tag{G32.2}
\]
Then
\[
 E_h=C_h\mathbb Q(\sqrt{-q},\sqrt{-d})
                                                               \tag{G32.3}
\]
admits every field of the bare source \texttt{Host}, with
$\mathfrak c_K=\mathcal O_{\mathbb Q(\sqrt{-d})}$.
Its exact numerical data include
\[
 \begin{split}
 [E_h:\mathbb Q]&=2^{h+3}=4n,\\
 \operatorname{rd}(E_h)
   &=2^{h+1-2^{-h}}\sqrt{3qd},\qquad
 \mu(E_h)=\mu_6.
 \end{split}                                                   \tag{G32.4}
\]
The source half-different gives the same principal identity
\[
 AA=\mathfrak b^{-1},\quad \xi_0=-1/\sqrt{-d},\qquad
        (\xi_0)AA\overline{AA}\mathfrak D_{E_h/\mathbb Q}=1.
                                                               \tag{G32.5}
\]

Let $\psi_{-3}$ be the nontrivial Dirichlet character of
conductor three.  This field satisfies the literal definition
\[
             \texttt{IsChiField}(\psi_{-3},h,q,d,E_h).
                                                               \tag{G32.6}
\]
Consequently any separately supplied instance
$\texttt{VAt}(\psi_{-3},\epsilon,401,q,d,h,x)$ applies
directly to the actual strict \emph{untwisted} groups over
$E_h$, with the exact cutoff $\lfloor x^{401}\rfloor$
and budget $\sqrt{\exp(x^\epsilon)}$.
No simultaneous-character window or quadratic descent is
needed for this particular application.  The proposition
neither proves \texttt{VWin} from A nor substitutes a
cubic Artin coefficient into its pure-Tate definition.
\end{proposition}
\begin{proof}
\proofstep{1. The real cyclotomic factor and the frame.}
Write $N=2^{h+1}$ and $z=\zeta_{2^{h+2}}$.
The integral basis $1,z,\ldots,z^{N-1}$ has minimal
polynomial $X^N+1$.  Comparing its coefficients under
conjugation shows
\[
 \mathcal O_{R_h}=\mathbb Z[z+z^{-1}].
\]
Indeed the conjugation-fixed integral basis is $1$ and
$z^j+z^{-j}$ for $1\leq j<N/2$; these are monic integral
polynomials of successive degrees in $z+z^{-1}$.
Also $\mathcal O_{\mathbb Q(z)}$ is free on $1,z$ over
this real ring, by the relation
$z^2-(z+z^{-1})z+1=0$.
Its relative discriminant is generated by
$(z-z^{-1})^2$.  Its absolute norm from $R_h$ is four:
$z^2$ is a primitive $2^{h+1}$-st root and
$|N_{\mathbb Q(z)/\mathbb Q}(1-z^2)|=\Phi_{2^{h+1}}(1)^2=4$.
The derivative of $X^N+1$ gives the full cyclotomic
absolute discriminant $2^{(h+1)2^{h+1}}$.
Transitivity therefore gives
\[
 |\operatorname{disc}R_h|=2^{(h+1)2^h-1},
 \qquad \operatorname{rd}(R_h)=2^{h+1-2^{-h}}.
                                                               \tag{G32.7}
\]
The integral cyclotomic basis used here is the standard
one in Milne, \emph{Algebraic Number Theory},
Theorem~6.4 and Remark~6.6, printed pp.~99--101.
The relative calculation above proves the precise
maximal-real exponent needed here.

The groups of $R_h$ and $K$ are disjoint by ramification.
The image of $5$ generates the cyclic real two-power
quotient, with order $2^h$, exactly as in AM-H17.
Its chosen lift acts trivially on $K$, so the kernel
condition on its cyclic subgroup holds.  This proves
the frame data, and $o^{64}=2^{64}\leq2^h$.
The displayed example of ranks satisfies all inequalities
in \texttt{Params.Ranks}.

\proofstep{2. Field coordinates and the local clauses.}
The four factors $R_h,K,K_q,K_d$ have disjoint conductor
supports $2,3,q,d$.  They are linearly disjoint, and their
integral bases tensor because the discriminants are
coprime.  Thus the degree and root discriminant in
(G32.4) follow from (G32.7) and the quadratic discriminants
$-3,-q,-d$.  These quadratic discriminants have the stated
form because $q,d\equiv3\pmod4$.

Let $c_C$ be conjugation on $C_h$; it is trivial on $R_h$
and nontrivial on $K$.  In physical coordinates use
\[
 (a,\epsilon_q,\epsilon_d)\longmapsto
 \bigl((\epsilon_q,\epsilon_d),
        a c_C^{(1-\epsilon_q)/2}\bigr).              \tag{G32.8}
\]
It takes full conjugation to $((-1,-1),1)$ and pure
$d$-inertia to $((1,-1),1)$.  Since $c_C$ projects to
one in the real two-power factor, it preserves the exact
\texttt{art\_h} condition.  Cyclotomic restriction from
conductor $Q_Sqd$ supplies the Artin map, its surjectivity,
the two quadratic coordinate characters, and the same
inertia, decomposition and Frobenius dictionary as AM-H17.
This dictionary includes the ramified primes: one removes
the prime-power part of the conductor before imposing the
congruence describing inertia, decomposition or Frobenius.
The exact declarations are \texttt{HostRaw} and
\texttt{HostDer} in
\texttt{ConjEMemHostFIff.lean}.

The conductor is exact at three by its nontrivial quadratic
factor.  At two the residue $1+2^{h+1}$ modulo
$2^{h+2}$ is neither $1$ nor $-1$ and acts nontrivially
on $R_h$; extending it by one at the other factors proves
the literal \texttt{QS\_exact} condition.  The clauses
$2^{h+2}\mid Q_S$, $(Q_S,qd)=1$ and $Q_S\leq n^2$
hold, and the only odd prime in $Q_S$ is three, with
$4\cdot3\leq B$.

At $2,3,q$, the decomposition group fixes $K_d$.
For two this follows from $d\equiv7\pmod8$; for three
and $q$ it follows from (G32.2).  Hence these groups
preserve the source CM type and exclude conjugation.
Quadratic reciprocity gives $(-q/d)=-1$:
$(d/q)=(-1/q)=-1$, hence $(q/d)=1$, and $(-1/d)=-1$.
If conjugation lay in decomposition at $d$, its restriction
to $K_q$ would force the Frobenius power to be odd.
On the real subfield of $\mathbb Q(\zeta_{16})$ this
odd power acts by the class of $7$ modulo $\{\pm1\}$,
which is nontrivial; conjugation acts trivially there.
This contradiction excludes it at $d$ as well.
Thus $E_h/E_h^+$ is unramified at every finite place,
and all the \texttt{decomp\_typeI}, \texttt{decompConj},
\texttt{unramEF}, \texttt{eta}, \texttt{ordinary} and
\texttt{splitK} clauses follow by the identical group
interpretations used in AM-H17.  The type in these
clauses is still the source type induced from $K_d$.

The displayed discriminant proves all its support clauses.
The required lower root-discriminant bound follows by
dropping $\sqrt{3q}$ from (G32.4); its upper bound follows
from
$2^{-2^{-h}}\sqrt3\leq3$.
Finally $\operatorname{rd}(E_h)^2\geq4n$, since already
$2h+2-2^{1-h}\geq h+3$ for $h\geq2$.
These check the three inequalities in \texttt{HostDer}.

\proofstep{3. The integral data and the roots of unity.}
The relative different of $E_h/K_d$ is supported at
$2,3,q$, all split in $K_d$.  At each such prime choose
one of its two primes in $K_d$ and put in $\mathfrak b$
all factors of the relative different above that choice,
with their full exponents.  Conjugation exchanges the
supports.  Thus
\[
 \mathfrak b\overline{\mathfrak b}=\mathfrak D_{E_h/K_d},
 \quad(\mathfrak b,\overline{\mathfrak b})=1,
 \quad \sigma\mathfrak b=\mathfrak b
                  \quad(\sigma\in\operatorname{Gal}(E_h/K_d)).
\]
The trace-dual ideal is exactly $\mathfrak b^{-1}$,
which verifies the source unimodularity clause for every
ideal with this factorization.  Transitivity of differents
and $\mathfrak D_{K_d/\mathbb Q}=(\sqrt{-d})$ give
(G32.5).

The unique quadratic subfield of $R_h$ is $\mathbb Q(\sqrt2)$.
Hence all quadratic subfields of $E_h$ have squareclasses
in the span of $2,-3,-q,-d$.  The squareclass $-1$ is
not in that span: valuations at $3,q,d$ first force their
three exponents to be zero, and neither $1$ nor $2$ is
$-1$ modulo rational squares.  In particular $i\notin E_h$.
A higher two-power root of unity is therefore impossible.
At three the exact ramification index is two, excluding
$\mu_9$; at $q$ and $d$ it is two, excluding their
primitive roots as $q,d>3$.  Every other odd prime is
unramified.  This proves $\mu(E_h)=\mu_6$.
The norm of each such root from $K$ to $\mathbb Q$ is
one, and its norm from $E_h$ to $K_d$ is a power of that
norm.  This verifies \texttt{norm\_mu}.

Take $\mathfrak c_K=\mathcal O_{K_d}$ and four distinct
level primes congruent to one modulo $15Q_Sqd$.
The same choices as in AM-H17 verify primality,
coprimality, membership and the two special-prime
congruences.  The source embedding is chosen positive
on $\sqrt{-d}/i$.  The bare choices
$TT=T_2=\varnothing$, $u_0=0$, $L_w=2$, $C_{T_2}=0$
verify the remaining stored fields and their vacuous
range conditions.  No later size lower bound on these
sets is asserted.  Thus all fields of \texttt{Host}
have been supplied.  As in AC-G30, the separately defined
\texttt{CoprimeFifteen} property fails because $3\mid Q_S$.

For completeness, primes satisfying (G32.2) exist without
any window assertion: choose $q>B$ in its residue class,
then choose $d$ in the coprime CRT class given there,
using Dirichlet's theorem.  The choice in the source
\texttt{WinPair} has these same residue conditions for
all odd primes below its stipulated splitting bound;
its much stronger size and window assertions are separate.

\proofstep{4. The literal cyclic field.}
Embed $E_h$ in the cyclotomic field of conductor $Q_Sqd$.
An automorphism given by a unit residue $a$ fixes it
exactly when
\[
 a\equiv\pm1\pmod {2^{h+2}},\quad
 \psi_{-3}(a)=1,\quad (a/q)=(a/d)=1.
\]
These are precisely the four conditions in
\texttt{ExistsPageN.IsChiField}, including the ambient
conductor divisibilities.  This proves (G32.6).
The assertion about \texttt{VAt} is its direct definition
applied to this field, not a new theorem supplying its
hypothesis.  Its \texttt{LData} remains the strict Tate
coefficient of \texttt{GaloisCohomology/Statement.lean}.
\end{proof}

\begin{proposition}[AC-G32-signature: the exact fixed-lattice CM-type gate]
\label{AC-G32-signature}
The maximal real subfield of (G32.3) is
\[
 E_h^+=R_h\mathbb Q(\sqrt{3q},\sqrt{3d}).             \tag{G32.9}
\]
There is a purely imaginary trace coefficient $\beta$
with $(\beta)=(\xi_0)$ making the unchanged lattice
$AA$ in (G32.5) positive for the type induced from $K$
if and only if there is a unit
$\kappa\in\mathcal O_{E_h^+}^{\times}$ satisfying
\[
 \operatorname{sign}\tau(\kappa)
       =\operatorname{sign}\tau(\sqrt{3d})
             \quad\text{for every real embedding }\tau.
                                                               \tag{G32.10}
\]
Here the imaginary square roots and the positive root of
$3d$ are chosen compatibly.  Such a unit cannot lie in
the quadratic subfield $\mathbb Q(\sqrt{3d})$.
This last obstruction alone makes no assertion about
existence in the larger field (G32.9).
\end{proposition}
\begin{proof}
Complex conjugation changes the signs of all three
imaginary quadratic roots and fixes $R_h$.  Its fixed
field is generated by their pairwise products, giving
(G32.9).

If $\beta$ has the stated ideal and is purely imaginary,
then $\kappa=\xi_0/\beta$ is a real integral unit;
conversely every real unit gives such a $\beta$.
Put $\beta_3=-1/\sqrt{-3}$, positive imaginary for the
chosen type on $K$.  Then
\[
 \frac{\beta}{\beta_3}
       =\frac{\sqrt{-3}}{\sqrt{-d}\,\kappa}
       =\frac{\sqrt{3d}}{d\,\kappa}.                 \tag{G32.11}
\]
The right-hand side is real.  It is positive at every
embedding exactly when (G32.10) holds.  This is precisely
the positivity condition for the induced type.
The equality of ideals leaves (G32.5) unchanged, so
when the condition holds the corresponding alternating
trace form is integral and principal.  This verifies
both directions of the stated fixed-lattice equivalence.

Now $d\equiv7\pmod {16}$ and $d\equiv-1\pmod3$.
The squarefree number $3d$ is $1$ modulo four, so a
unit of its real quadratic field is
$(a+b\sqrt{3d})/2$ with integers $a,b$ of matching parity.
Opposite signs at its two embeddings would force norm
$-1$, hence
\[
                         a^2-3db^2=-4.
\]
Reducing modulo three gives $a^2=2$, impossible in
$\mathbb F_3$.  This proves the quadratic obstruction.
The extension from that quadratic field to (G32.9) has
even degree; a norm from it squares this signature.
Thus the same argument cannot be used to exclude a
unit in the larger field.
\end{proof}

\paragraph{What is and is not repaired.}
The smaller construction supplies a cubic-compatible bare Host
and one exact cyclic-character field.  It removes the specific
multi-character field-definition obstruction of (G30.7) for
this choice of host.  It does not retain the inherited
$\mathbb Q(\zeta_7)$ Jacobian and its specified sign unit:
neither $\zeta_7$ nor $i$ lies in the new field.
The AM-H16 units live in explicitly larger real subfields;
weak approximation on signs does not replace them by units
in (G32.9).  AC-G33 below constructs the required unit in (G32.10)
with an explicit logarithmic height bound.  The equivalence
here isolates exactly the condition that construction meets.
The cross-block coefficient tie of AC-G29--G30 is
still needed for the cubic Artin character.  The source
\texttt{VVle} hypothesis obtained from a supplied
\texttt{VWin} continues to bound pure-Tate strict groups
over this varying field, not those mixed coefficients.

\paragraph{Algebraic-geometric dictionary.}
The finite cyclic character in (G32.6) corresponds to a finite
etale covering after the conductor primes are inverted, in the
Galois-category dictionary of SGA~1, Expos\'e~V, \S\S4 and 7,
already used in AC-K19.  This records the finite coefficient
and its field; the principal CM lattice and the later theta
comparison are additional structures.  No new archive cote
is asserted here.
